\documentclass[10pt,a4paper]{amsart}
\usepackage[margin=19mm]{geometry}
\usepackage{amsmath,amssymb,mathtools}
\usepackage[scr=rsfs,cal=euler]{mathalfa}
\usepackage{xfrac}
\usepackage[unicode,hidelinks,bookmarks=false]{hyperref}
\newcommand{\ie}{i.e.\ }
\newcommand{\cf}{cf.\ }
\newcommand{\resp}{respectively\ }
\newcommand{\Subsection}{\subsection}
\providecommand{\nobreakdash}{\nobreak\hbox{-}}
\setlength{\emergencystretch}{2em}
\multlinegap0pt
\usepackage{bm}
\usepackage{bbm}
\usepackage{dsfont}
\usepackage{xspace}
\usepackage{cancel}
\usepackage{mathtools}
\usepackage{cleveref}
\usepackage{amsthm}
\makeatletter
\@ifundefined{prop}{\newtheorem{prop}{Prop}}{}
\@ifundefined{claim}{\newtheorem{claim}[prop]{Claim}}{}
\@ifundefined{cor}{\newtheorem{cor}[prop]{Corollary}}{}
\@ifundefined{lemma}{\newtheorem{lemma}[prop]{Lemma}}{}
\makeatother
\newcommand{\Z}{\ensuremath{\mathbf{Z}}}
\newcommand{\md}[1]{\ensuremath{(\operatorname{mod}\, #1)}}
\newcommand{\ru}[1]{e_{g-1}(#1)}
\renewcommand{\geq}{\geqslant}
\renewcommand{\leq}{\leqslant}
\title[Niven numbers are an asymptotic basis of order 3]{Niven numbers are an asymptotic basis of order 3}
\author{Kate Thomas}
\date{Source: arXiv:2511.02790v1 (CC BY 4.0)}
\begin{document}
\maketitle

\noindent\emph{Source and licence.} Niven numbers are an asymptotic basis of order 3. Version arXiv:2511.02790v1 (CC BY 4.0), distributed under the Creative Commons Attribution 4.0 International licence, as recorded in the arXiv licence metadata for this exact version and shown on the abstract page https://arxiv.org/abs/2511.02790v1. Full rights notice: licenses/arxiv-2511.02790-CC-BY-4.0.txt.\par\smallskip

\noindent\textit{Editorial scope.} Counted unit: the Lemma whose source label is \texttt{psi lemma} in the article (source0.tex lines 361--364, source label \texttt{psi lemma}), delivered together with the article's own complete original proof of it (source0.tex lines 365--521, one \texttt{proof} environment of 157 lines). No mathematics is written, repaired or re-derived: statement and proof are the article's own text line for line. Required context retained uncounted: Local limit theorem application version (lines 319--324); psi defn (lines 335--337); j range remark (lines 339--341). Every definition, display and label of those lines is retained unchanged. Omitted from this package and not redistributed: 1--318, 325--334, 338--338, 342--360, 522--1208. Nothing inside a retained excerpt is omitted or altered: every retained body is delivered exactly as the article's lines are written, so no omission receipt of any kind accompanies this package. Citations: the retained excerpts carry no citation command at all, so no bibliography entry of the article is transcribed and no reference is redistributed. Reference adaptation: none was needed and none was made; every reference inside the delivered unit resolves inside it, so no unresolved reference or citation is produced. Printed slips kept literal: no printed word, symbol, hypothesis or number of the counted statement or of its proof was altered. Layout and class: the article's own class is \texttt{amsart} with packages including graphicx, todonotes, amsfonts, amsthm, amsmath, parskip, mathrsfs, amssymb, xtab, dsfont, mathtools, natbib, csquotes, url; the wrapper is the lane's pinned \texttt{amsart} editorial wrapper at 10pt on A4, which restates the theorem-like environments the retained text needs and adds the article's own macro definitions that the retained text needs (\texttt{\textbackslash Z}, \texttt{\textbackslash geq}, \texttt{\textbackslash leq}, \texttt{\textbackslash md}, \texttt{\textbackslash ru}), with extra wrapper packages (bm, bbm, dsfont, xspace, cancel, mathtools, cleveref). The delivered numbering is the unit's own. Assets: no retained span contains a figure environment, a graphics-inclusion command, a PDF-page inclusion, a table or a TikZ picture; no image file is copied into this package, the package declares no asset dependency and no image file is redistributed with it. Licence: version arXiv:2511.02790v1 of this article is distributed under the Creative Commons Attribution 4.0 International licence, as recorded in the arXiv licence metadata for this exact version and shown on the abstract page https://arxiv.org/abs/2511.02790v1. A full rights notice is delivered at licenses/arxiv-2511.02790-CC-BY-4.0.txt. Characters: every retained excerpt is pure ASCII, so the delivered engine, XeLaTeX with its default Latin Modern text font, reports no missing character.
\begin{cor}\label{Local limit theorem application version}
    For $T,\nu$ and $P(T,\nu)$ as defined above,
    \[P(T,\nu)=\frac{e^{-x^{2}/2}}{\sqrt{2\pi\sigma^{2}T}}+O_{g}(T^{-3/2}),\]
    where $x=\nu/\sqrt{\sigma^{2}T}$, and $\sigma^{2}=(g^{2}-1)/12$ is the variance of $X$. In particular, if $|x|<1/2$,
    \[P(T,\nu)=\frac{1}{\sqrt{2\pi\sigma^{2}T}}+O_{g}(\max (x^{2}T^{-1/2},T^{-3/2})).\]
\end{cor}

\begin{equation}\label{psi defn}
    \Psi(\mathbf{a};\nu)\coloneq \sum_{\substack{j_{1},\ldots, j_{g-2}\\ \sum_{i=0}^{g-2} j_{i}=\nu}}\prod_{s=0}^{g-2}\ru{sj_{s}}P(a_{s},j_{s}).
\end{equation}

\begin{equation}\label{j range remark}
|j_{s}|\leq a_{s}(g-1)/2 \textrm{ for }s=1,\ldots, g-2.
\end{equation}

\begin{lemma}\label{psi lemma}
    Suppose that $a_{0},\ldots, a_{g-2}\geq 0$ be integers with $\mathbf{a}=(a_{0},\ldots, a_{g-2})$ and let $\nu\in \Z$. Suppose further that $|\nu|\leq Ca_{0}^{1/4}$ for some $C> 0$ and $a_{0}\geq a_{s}$ for $s=1,\ldots, g-2$. Then 
    \[\Psi(\mathbf{a};\nu)=\Psi(\mathbf{a};0)+O_{C,g}(\nu^{2}a_{0}^{-3/2}).\]
\end{lemma}

\begin{proof}
From the definition \cref{psi defn}, we see that
\begin{equation}\label{psi defn restated}
    \Psi(\mathbf{a};\nu)=\sum_{j_{1},\ldots, j_{g-2}}P(a_{0},\nu-\sum_{r=1}^{g-1}j_{r})\prod_{s=1}^{g-2}\ru{sj_{s}}P(a_{s},j_{s}).
\end{equation}
    First, we use the local limit theorem to estimate the term $P(a_{0},\nu-\sum_{r=1}^{g-2}j_{r})$. Let $J\coloneqq \sum_{r=1}^{g-2}j_{r}$. Note that $|J|< g^{2}\max a_{i}$ from \cref{j range remark}. From \cref{Local limit theorem application version}, we have 
    \begin{equation}\label{psi 0}
        P(a_{0},\nu-J)=(2\pi\sigma^{2}a_{0})^{-1/2}e^{-(\nu-J)^{2}/2\sigma^{2}a_{0}}+O_{g}(a_{0}^{-3/2}).
    \end{equation}
    Substituting this into \cref{psi defn restated} gives
    \begin{equation}\label{psi 1}
        \Psi(\mathbf{a};\nu)=(\sigma\sqrt{2\pi a_{0}})^{-1}e^{-\nu^{2}/2\sigma^{2}a_{0}}\sum_{j_{1},\ldots, j_{g-2}}e^{(2\nu J-J^{2})/2\sigma^{2}a_{0}}\prod_{s=1}^{g-2}\ru{sj_{s}}P(a_{s},j_{s})+O_{g}(a_{0}^{-3/2}),
    \end{equation}
    where the error term in \cref{psi 1} comes from that of \cref{psi 0}, and the fact that
    \begin{equation}\label{product of probs}
        \sum_{j_{1},\ldots, j_{g-2}}\prod_{s=1}^{g-2}P(a_{s},j_{s})=1.
    \end{equation}
By expanding the term $e^{-\nu^{2}/2\sigma^{2}a_{0}}=1+O_{C,g}(\nu^{2}a_{0}^{-1})$ in \cref{psi 1}, we will show that:
    \begin{claim}\label{psi 2 claim}
        \begin{equation}\label{psi 2 equation}
            \Psi(\mathbf{a};\nu)=(\sigma\sqrt{2\pi a_{0}})^{-1}\sum_{j_{1},\ldots, j_{g-2}}e^{(2\nu J-J^{2})/2\sigma^{2}a_{0}}\prod_{s=1}^{g-2}\ru{sj_{s}}P(a_{s},j_{s})+O_{C,g}(\nu^{2}a_{0}^{-3/2}).
        \end{equation}
    \end{claim} To establish the claim, we first show that $e^{(2\nu J-J^{2})/2\sigma^{2}a_{0}}\ll_{C,g}1$ by considering the ranges $|\nu|\leq |J|/2$ and $|\nu|>|J|/2$ separately. In the former case, when $|\nu|\leq |J|/2$,
    \[e^{(2\nu J-J^{2})/2\sigma^{2}a_{0}}\leq e^{(2|\nu J|-J^{2})/2\sigma^{2}a_{0}}\leq 1.\]
    In the latter case, when $|\nu|>|J|/2$, we use the assumption that $|\nu|\leq Ca_{0}^{1/4}$ to obtain
    \[e^{(2\nu J-J^{2})/2\sigma^{2}a_{0}}\leq e^{2\nu^{2}/\sigma^{2}a_{0}}\leq e^{2C^{2}/\sigma^{2}a_{0}^{1/2}}\ll_{C,g} 1.\]
    Thus using that $e^{(2\nu J-J^{2})/2\sigma^{2}a_{0}}\ll_{C,g}1$ and \cref{product of probs} gives
    \[\sum_{j_{1},\ldots, j_{g-2}}e^{(2\nu J-J^{2})/2\sigma^{2}a_{0}}\prod_{s=1}^{g-2}P(a_{s},j_{s})\ll_{C,g}1.\]
    This concludes the proof of \cref{psi 2 claim}.

    Now we discard the terms with $|J|\geq \sigma^{2}a_{0}/2|\nu|$ from \cref{psi 2 equation} (if there are any), showing that the contribution from these terms is negligible. Indeed, as $\sigma^{2}a_{0}/2|\nu|\leq |J|\leq g^{2}\max a_{s}$, and $\max a_{s}\leq a_{0}$,
    \begin{equation}\label{psi exp bound}
        e^{(2\nu J-J^{2})/2\sigma^{2}a_{0}}\leq e^{g^{2}|\nu|\max a_{s}/\sigma^{2}a_{0} -\sigma^{2}a_{0}/8\nu^{2}}\leq e^{Cg^{2}a_{0}^{1/4}/\sigma^{2}-\sigma^{2}a_{0}^{1/2}/8C^{2}}\ll_{C,g} a_{0}^{-10}.
    \end{equation}
    This additionally uses the assumption that $|\nu|\leq Ca_{0}^{1/4}$ for some $C>0$. Bounding the term $|\ru{sj_{s}}P(a_{s},j_{s})|\leq 1$ for all $s, j_{s}$ and using \cref{psi exp bound}, we have
    \[\Big | \sum_{\substack{j_{1},\ldots, j_{g-2}\\ \sigma^{2}a_{0}/2|\nu|\leq|J|\leq g^{2}a_{0}}} e^{(2\nu J-J^{2})/2\sigma^{2}a_{0}}\prod_{s=1}^{g-2}\ru{sj_{s}}P(a_{s},j_{s})\Big |\ll_{C,g} a_{0}^{-9},\]
    hence the contribution to \cref{psi 2 equation} from $(j_{1},\ldots, j_{g-2})$ such that $|J|$ is large can be absorbed into the overall error term of $O_{C,g}(\nu^{2}a_{0}^{-3/2})$. It remains to estimate the contribution to \cref{psi 2 equation} from $(j_{1},\ldots, j_{g-2}) $ with $|J|\leq \sigma^{2}a_{0}/2|\nu|$. To this end, we expand the term $e^{\nu J/\sigma^{2}a_{0}}$ in \cref{psi 2 equation}, giving
    \begin{equation}\label{psi proof eq}
        \Psi(\mathbf{a};\nu)=\frac{1}{\sigma\sqrt{2\pi a_{0}}}\sum_{\substack{j_{1},\ldots, j_{g-2}\\|J|\leq \sigma^{2}a_{0}/2|\nu|}}e^{-J^{2}/2\sigma^{2}a_{0}}\Big (1+\frac{\nu J}{\sigma^{2}a_{0}}+O_{g}\Big(\frac{\nu^{2}J^{2}}{a_{0}^{2}}\Big)\Big )\prod_{s=1}^{g-2}\ru{sj_{s}}P(a_{s},j_{s})+O_{C,g}(\nu^{2}a_{0}^{-3/2}).
    \end{equation}
    First, we show that
    \begin{equation}\label{psi main term}
        \frac{1}{\sigma\sqrt{2\pi a_{0}}}\sum_{\substack{j_{1},\ldots, j_{g-2}\\|J|\leq \sigma^{2}a_{0}/2|\nu|}}e^{-J^{2}/2\sigma^{2}a_{0}}\prod_{s=1}^{g-2}\ru{sj_{s}}P(a_{s},j_{s})=\Psi(\mathbf{a};0)+O_{C,g}(a_{0}^{-3/2}).
    \end{equation}
    To do so, note we can undo the truncation on the range of summation in \cref{psi main term}. If the contribution from the range $\sigma^{2}a_{0}/2|\nu|\leq |J|\leq g^{2}a_{0}$ is non-zero, then in this range:
    \[e^{-J^{2}/2\sigma^{2}a_{0}}\leq e^{-\sigma^{2}a_{0}/8\nu^{2}}\leq e^{-\sigma^{2}a_{0}^{1/2}/8C^{2}}\ll_{C,g}a_{0}^{-10}.\]
    Using this and \cref{product of probs} we have
    \begin{equation}\label{range truncation}
        \sum_{\substack{j_{1},\ldots, j_{g-2}\\|J|> \sigma^{2}a_{0}/2|\nu|}}e^{-J^{2}/2\sigma^{2}a_{0}}\prod_{s=1}^{g-2}\ru{sj_{s}}P(a_{s},j_{s})\ll_{C,g} a_{0}^{-10}.
    \end{equation}
  
    We now work in the full range of summation for $(j_{1},\ldots, j_{g-2})$. We apply \cref{Local limit theorem application version} to note that $e^{-J^{2}/2\sigma^{2}a_{0}}/\sigma\sqrt{2\pi a_{0}}=P(a_{0},-J)+O_{g}(a_{0}^{-3/2})$, hence
    \begin{align*}
         \frac{1}{\sigma\sqrt{2\pi a_{0}}}&\sum_{j_{1},\ldots, j_{g-2}}e^{-J^{2}/2\sigma^{2}a_{0}}\prod_{s=1}^{g-2}\ru{sj_{s}}P(a_{s},j_{s})\\&=\sum_{j_{1},\ldots, j_{g-2}}P(a_{0},-J)\prod_{s=1}^{g-2}\ru{sj_{s}}P(a_{s},j_{s})+O_{C,g}(a_{0}^{-3/2})\\
         &=\Psi(\mathbf{a};0)+O_{C,g}(a_{0}^{-3/2}),
    \end{align*}
    using the definition in \cref{psi defn} to obtain the second equality. This establishes \cref{psi main term}, giving the main term in \cref{psi lemma}.

   We now show that the remaining terms in \cref{psi proof eq} contribute only to the error term in \cref{psi lemma}. The $O_{g}(\nu^{2}J^{2}a_{0}^{-2})$ term within the summation over $(j_{1},\ldots, j_{g-2})$ in \cref{psi proof eq} can be absorbed into the error term $O_{C,g}(\nu^{2}a_{0}^{-3/2})$. Indeed, as $e^{-J^{2}/2\sigma^{2}a_{0}}J^{2}\leq 2\sigma^{2}a_{0}/e$, the contribution from this term is bounded as follows:
    \[\frac{\nu^{2}}{a_{0}^{5/2}}\sum_{\substack{j_{1},\ldots, j_{g-2}\\|J|\leq \sigma^{2}a_{0}/2|\nu|}}\prod_{s=1}^{g-2}P(a_{s},j_{s})e^{-J^{2}/2\sigma^{2}a_{0}}J^{2}\ll_{g} \frac{\nu^{2}}{a_{0}^{3/2}}\sum_{\substack{j_{1},\ldots, j_{g-2}\\|J|\leq \sigma^{2}a_{0}/2|\nu|}}\prod_{s=1}^{g-2}P(a_{s},j_{s})\ll_{g} \frac{\nu^{2}}{a_{0}^{3/2}},\]
    additionally using \cref{product of probs} in the final inequality. 

    We have shown that
    \begin{equation}\label{psi proof eq 2}
        \Psi(\mathbf{a};\nu)=\Psi(\mathbf{a};0)+\frac{\nu}{\sigma^{3}a_{0}^{3/2}\sqrt{2\pi}}\sum_{j_{1},\ldots, j_{g-2}}e^{-J^{2}/2\sigma^{2}a_{0}}J\prod_{s=1}^{g-2}\ru{sj_{s}}P(a_{s},j_{s})+O_{C,g}(\nu^{2}a_{0}^{-3/2}).
    \end{equation}
    Finally we show that \cref{psi proof eq 2} implies the statement of the lemma. It suffices to prove
    \begin{equation}\label{psi lemma final}
        \sum_{j_{1},\ldots, j_{g-2}}e^{-J^{2}/2\sigma^{2}a_{0}}J\prod_{s=1}^{g-2}\ru{sj_{s}}P(a_{s},j_{s})\ll_{C,g} 1.
    \end{equation}
Before proving \cref{psi lemma final}, we remark that the proof is essentially trivial in base 3. In this case, the equation on the left hand side of \cref{psi lemma final} equals
\[\sum_{j}(-1)^{j}je^{-j^{2}/2\sigma^{2}a_{0}}P(a_{1},j)=\sum_{j}g(j)=0,\]
as $g(j)\coloneq (-1)^{j}je^{-j^{2}/2\sigma^{2}a_{0}}P(a_{1},j)$ is an odd function. For $g\geq 4$, the proof of \cref{psi lemma final} is rather more involved, and we make use of the cancellation coming from the $\ru{sj_{s}}$ terms instead of the sign of $J$. Let
\begin{equation}\label{g defn}
        g(j_{1},\ldots, j_{g-2})\coloneq Je^{-J^{2}/2\sigma^{2}a_{0}}\prod_{s=1}^{g-2}P(a_{s},j_{s}).
    \end{equation}
We first work with the tuples $(j_{1},\ldots, j_{g-2})$ for which the following condition holds, in addition to the assumption throughout that $|j_{s}|\leq a_{s}(g-1)/2$. Suppose that $J=j_{1}+\ldots+j_{g-2}$ is such that
\begin{equation}\label{J tuples}
    \Big|\frac{2Jr+r^{2}}{2\sigma^{2}a_{0}}\Big |\leq 1/2\textrm{ for all  }r\in \{0,\ldots, g-2\}.
\end{equation}
For such $(j_{1},\ldots, j_{g-2})$, the function $g(j_{1},\ldots, j_{g-2})$ doesn't vary too much when incrementing $j_{1}$ by a small amount, as shown in the following claim.
\begin{claim}\label{g claim}
    For $(j_{1},\ldots, j_{g-2})$ such that \cref{J tuples} holds, and for all $r=0,\ldots, g-2$,
    \[|g(j_{1},\ldots, j_{g-2})-g(j_{1}+r,j_{2},\ldots, j_{g-2})|\ll_{g} e^{-J^{2}/2\sigma^{2}a_{0}}\Big (\frac{J^{2}}{a_{0}}(P(a_{1},j_{1})+a_{1}^{-3/2})+P(a_{1},j_{1})\Big )\prod_{s=2}^{g-2}P(a_{s},j_{s}).\]
\end{claim}

\begin{proof}[Proof of \cref{g claim}]
First we note that the claim is trivial when $J=0$. In this case, $(j_{1},\ldots, j_{g-2})=(0,\ldots, 0)$ and $g(0,\ldots, 0)=0$, and it follows from \cref{g defn} that $g(r,0,\ldots, 0)\ll_{g} \prod_{s=1}^{g-1}P(a_{s},j_{s})$. Thus from now on, we assume that $|J|\geq 1/2$. From \cref{Local limit theorem application version},
\begin{equation}\label{prob +r roughly const}
    |P(a_{1},j_{1})-P(a_{1},j_{1}+r)|\ll_{g}a_{1}^{-1/2}|e^{-a_{1}^{2}/2\sigma^{2}a_{1}}-e^{-(a_{1}+r)^{2}/2\sigma^{2}a_{1}}|+O_{g}(a_{1}^{-3/2}) \ll_{g} a_{1}^{-3/2}.
\end{equation}
Here, we also use that the function $e^{-x^{2}}$ is Lipschitz to obtain the final inequality. We also have, under the assumptions of \cref{J term suff cancellation claim}, 
\begin{equation}\label{exp Jr term}
    \exp\Big (-\frac{2Jr+r^{2}}{2\sigma^{2}a_{0}} \Big )=1+O_{g}\Big (\frac{J}{a_{0}}\Big ).
\end{equation}
From the definition in \cref{g defn}, 
\begin{align*}
    & |g(j_{1},\ldots, j_{g-2})-g(j_{1}+r,j_{2},\ldots, j_{g-2})|\\
     &=e^{-J^{2}/2\sigma^{2}a_{0}}\prod_{s=2}^{g-2}P(a_{s},j_{s})|JP(a_{1},j_{1})-(J+r)e^{-(2Jr+r^{2})/2\sigma^{2}a_{0}}P(a_{1},j_{1}+r)|.
\end{align*}
Expanding the term $e^{-(2Jr+r^{2})/2\sigma^{2}a_{0}}$ using \cref{exp Jr term}, and using \cref{prob +r roughly const} to estimate $P(a_{1}+r,j_{1})$, we obtain,
\begin{align*}
    |JP(a_{1},j_{1})-(J+r)e^{-(2Jr+r^{2})/2\sigma^{2}a_{0}}&P(a_{1},j_{1}+r)|\\&=|JP(a_{1},j_{1})-(J+r)(1+O_{g}(J/a_{0}))(P(a_{1},j_{1})+O_{g}(a_{1}^{-3/2})))|\\
    &\ll_{g}\frac{J^{2}}{a_{0}}(P(a_{1},j_{1})+a_{1}^{-3/2})+P(a_{1},j_{1}).
\end{align*}
To simplify the final expression, we have used that $r\ll_{g} 1$, and that for any tuple $(j_{1},\ldots, j_{g-2})$, $J^{2}\geq |J|/2$. 
\end{proof}
We use \cref{g claim} to prove the following:
\begin{equation}\label{J term suff cancellation claim}
     \sum_{j_{1},\ldots, j_{g-2}}e^{-J^{2}/2\sigma^{2}a_{0}}J\prod_{s=1}^{g-2}\ru{sj_{s}}P(a_{s},j_{s})\ll_{g} 1.
\end{equation} 
We have
\begin{equation*}
    \sum_{j_{1},\ldots, j_{g-2}}e^{-J^{2}/2\sigma^{2}a_{0}}J\prod_{s=1}^{g-2}\ru{sj_{s}}P(a_{s},j_{s})=\sum_{j_{1},\ldots, j_{g-2}}g(j_{1},\ldots, j_{g-2})\prod_{s=1}^{g-2}\ru{sj_{s}}
\end{equation*}
To get the cancellation required, we use \cref{g claim} to assert that $g(j_{1}+r,\ldots, j_{g-2})$ is essentially constant as $r$ varies in $\{0,\ldots, g-2\}$, which allows us to get cancellation from the $\prod_{s=1}^{g-2}\ru{sj_{s}}$ term. Our aim is to split the range of summation into sets where $j_{1}$ has a fixed congruence modulo $g-1$. Note that if $g$ is even and $a_{1}$ odd, the range of $j_{1}$ is not contained in the integers, as $j_{1} \in \{0,\ldots, a_{1}(g-1)\}-a_{1}(g-1)/2$. In this case, $j_{1}+1/2\subset \Z$, so we can run the following argument by multiplying through by a factor of $\ru{1/2}$. If $a_{1}(g-1)/2\in \Z$, let $x\coloneq -a_{1}(g-1)/2$, otherwise let $x\coloneq -a_{1}(g-1)/2-1/2$. Note that from \cref{j range remark}, the range of $j_{1}$ is a multiple of $g-1$. Thus splitting the range of $j_{1}$ (compensating by a factor of $\ru{1/2}$ if necessary), 
\begin{align}
    \sum_{j_{1},\ldots, j_{g-2}}g(j_{1},\ldots, j_{g-2})\prod_{s=1}^{g-2}\ru{sj_{s}}=\sum_{\substack{j_{1},\ldots, j_{g-2}\\ j_{1}\equiv x \md{g-1}}}\prod_{s=1}^{g-2}\ru{sj_{s}}\sum_{r=0}^{g-2}\ru{r}g(j_{1}+r,j_{2},\ldots, j_{g-2})\nonumber\\
    =\sum_{\substack{j_{1},\ldots, j_{g-2}\\ j_{1}\equiv x \md{g-1}}}\prod_{s=1}^{g-2}\ru{sj_{s}}\sum_{r=0}^{g-2}\ru{r}\Big ( g(j_{1},j_{2},\ldots, j_{g-2})+O_{g}(E(J,a_{0},a_{1}))\Big )\label{final g sum}
\end{align}
from \cref{J term suff cancellation claim}, where
\[E(J,a_{0},a_{1})=e^{-J^{2}/2\sigma^{2}a_{0}}\Big (\frac{J^{2}}{a_{0}}(P(a_{1},j_{1})+a_{1}^{-3/2})+P(a_{1},j_{1})\Big )\prod_{s=2}^{g-2}P(a_{s},j_{s}).\]
The first term in the sum over $r$ is zero:
\begin{align*}
    &\sum_{\substack{j_{1},\ldots, j_{g-2}\\ j_{1}\equiv x \md{g-1}}}\prod_{s=1}^{g-2}\ru{sj_{s}}\sum_{r=0}^{g-2}\ru{r}g(j_{1},j_{2},\ldots, j_{g-2})\\&=\sum_{\substack{j_{1},\ldots, j_{g-2}\\ j_{1}\equiv x \md{g-1}}}\prod_{s=1}^{g-2}\ru{sj_{s}}g(j_{1},j_{2},\ldots, j_{g-2})\sum_{r=0}^{g-2}\ru{r}=0.
\end{align*}
Finally, we show that the error term in \cref{final g sum} is $\ll_{g} 1$, that is, we show
\begin{equation}\label{final g sum error}
    \sum_{j_{1},\ldots, j_{g-2}}E(J,a_{0},a_{1})\ll_{g} 1.
\end{equation}
First note that from \cref{product of probs}, 
\begin{equation}\label{no J sq term}
    \sum_{j_{1},\ldots, j_{g-2}}e^{-J^{2}/2\sigma^{2}a_{0}}\prod_{s=1}^{g-2}P(a_{s},j_{s})\ll_{g} 1.
\end{equation}
and further using that $\sup J^{2}e^{-J^{2}/2\sigma^{2}a_{0}}\ll_{g} a_{0}$ gives
\begin{equation}\label{J sq 1}
    \frac{1}{a_{0}}\sum_{j_{1},\ldots, j_{g-2}}J^{2}e^{-J^{2}/2\sigma^{2}a_{0}}\prod_{s=1}^{g-2}P(a_{s},j_{s})\ll_{g} 1.
\end{equation}
    Moreover, a similar statement holds when replacing $P(a_{1},j_{1})$ by $a_{1}^{-3/2}$:
    \begin{align}
         \frac{1}{a_{0}a_{1}^{3/2}}\sum_{j_{1},\ldots, j_{g-2}}J^{2}e^{-J^{2}/2\sigma^{2}a_{0}}\prod_{s=2}^{g-2}P(a_{s},j_{s})&\ll_{g} \sum_{j_{1}}\frac{1}{a_{1}^{3/2}}\sum_{j_{2},\ldots, g-2}\prod_{s=2}^{g-2}P(a_{s},j_{s})\nonumber\\
         &=\sum_{j_{1}}\frac{1}{a_{1}^{3/2}}\ll_{g} \frac{1}{a_{1}^{1/2}},\label{J sq 2}
    \end{align}
    where the last line uses that $j_{1}$ ranges over $|j_{1}|\leq a_{1}(g-1)/2$.
    Combining \cref{no J sq term}, \cref{J sq 1} and \cref{J sq 2} gives \cref{final g sum error}.

Assuming \cref{J term suff cancellation claim}, it suffices to show that the contribution from $(j_{1},\ldots, j_{g-2})$ such that \cref{J tuples} doesn't hold is bounded. Suppose for $(j_{1},\ldots, j_{g-2})$ and $J=j_{1}+\ldots +j_{g-2}$,
\[\Big|\frac{2Jr+r^{2}}{2\sigma^{2}a_{0}}\Big |>\frac{1}{2}.\]
In particular, for such $J$, $|J|>\sigma^{2}a_{0}/2-(g-1)/2$, therefore
\begin{equation*}
    J^{2}e^{-J^{2}/2\sigma^{2}a_{0}}\ll_{g} a_{0}^{2}e^{-\sigma^{2}a_{0}/8}\ll_{g} a_{0}^{-10}.
\end{equation*}
Thus the contribution from $j_{1},\ldots, j_{g-2}$ where $|J|$ is large is
\begin{align*}
    \Big |\sum_{\substack{j_{1},\ldots, j_{g-2}\\ |J|\geq \sigma^{2}a_{0}/4}}J^{2}e^{-J^{2}/2\sigma^{2}a_{0}}\prod_{s=1}^{g-2}P(a_{s},j_{s})\Big |&\ll_{g} a_{0}^{-10}.\qedhere
\end{align*}
\end{proof}

\end{document}
